\chapter{Conclusion}
\label{ch:conclusion}

\section{Obtained Results}

\textbf{Template intent.} Keep this section high-level until final numbers are fixed.

\begin{itemize}
	\item Summarize what was built and evaluated.
	\item Restate outcome at the level of evidence (not hype): where ML improved, matched, or underperformed heuristic logic.
	\item Mention both decision quality and operational feasibility.
\end{itemize}

\section{Limitations}

\textbf{Template intent.} Be explicit and concrete.

\begin{itemize}
	\item Data limitations (coverage, label ambiguity, special race conditions).
	\item Method limitations (model class, assumptions, calibration drift).
	\item System limitations (latency budget, integration constraints, observability).
\end{itemize}

\section{Future Work}

\textbf{Template intent.} Propose realistic next steps tied to thesis findings.

\begin{itemize}
	\item Drift-aware adaptation and online recalibration.
	\item Richer decision spaces (for example compound recommendation jointly with timing).
	\item Tighter integration patterns and broader season-level validation.
	\item Training-serving skew mitigation when telemetry is streamed in parallel to Flink and ML inference.
\end{itemize}

Suggested references: \citep{gama2014drift, arora2024driftreview, bifet2007adwin, gomes2017arf, baenagarcia2006eddm}.

\subsection{Deployment Consistency Placeholder (Flink + ML)}

\textbf{Template intent.} Keep this subsection as a future-work marker until deployment experiments are run.

\begin{itemize}
	\item Define how online features should be generated at serving time to match training assumptions.
	\item Specify consistency checks between offline Flink-derived training features and live serving features.
	\item Reserve latency and fallback requirements for the final system design chapter pass.
\end{itemize}

\section{Overall Discussion}

\textbf{Template intent.} Close by positioning your thesis against the three strategy paradigms.

\begin{itemize}
	\item Optimization/planning paradigm.
	\item Reinforcement-learning paradigm.
	\item Supervised classification in streaming deployment.
\end{itemize}

Suggested references: \citep{carrasco2023optimization, quiroga2024optimization, thomas2025xrlf1, sasikumar2025pitstop, hettmann2024datatopodium, dinh2025gradient}.
